% Abstract and keywords. Every number is a \kn{} macro read from a results file.
\begin{abstract}
A Farmer Producer Organization (FPO) must decide where, when and how much of its members' produce is sold, and existing tools treat
forecasting, storage and logistics separately. KisanNiti connects them in five stages that a manager follows: collecting and cleaning
prices and farmer records, proposing a plan, reviewing it, telling each farmer in their language, and recording what was sold. A regime
test separates price-bound from market crops: for Paddy (Common) in Punjab, at least \kn{msp_share_min}\% of reported prices in every season
from \kn{season_first} to \kn{season_last} equal the minimum support price, and a rule predicting it had the lowest error (Rs~\kn{mae_msp_mean}
per quintal). For basmati, eight forecasting methods were compared under a protocol frozen before any test result existed; on the
held-out season none of \kn{inf_total} comparisons with the naive forecast at 1, 7 and 14 days found a method better after correcting for
multiple comparisons. An exact fair-allocation model raised the smallest share a farmer sells from \kn{allpol_p1_nearest_min_fulfilment_pct}\% under the
nearest-mandi rule to \kn{allpol_kisaniti_min_fulfilment_pct}\% in \kn{allpol_plans} synthetic plans at real prices, at Rs~\kn{allpol_kisaniti_price_of_fairness_rs_per_q} per
quintal. The system was built and tested end to end (\kn{sys_tests_total} automated tests, a \kn{sys_chain_steps}-step run on fictional data); no
real farmer data or user study is included.
\keywords{Farmer Producer Organizations \and Mandi price forecasting \and Minimum support price \and Fair allocation \and Integer
programming}
\end{abstract}
